% Generated draft: requires mathematical and source review.
\documentclass[11pt]{article}
\usepackage[margin=25mm]{geometry}
\usepackage{fontspec}
\setmainfont{Latin Modern Roman}
\usepackage{amsmath,amssymb,mathtools,unicode-math}
\setmathfont{Latin Modern Math}
\usepackage{enumitem,needspace,xurl,hyperref}
\setlength{\parindent}{0pt}
\setlength{\parskip}{5pt}
\setlength{\emergencystretch}{3em}
\newcommand{\N}{\mathbb{N}}
\newcommand{\Z}{\mathbb{Z}}
\newcommand{\Q}{\mathbb{Q}}
\newcommand{\R}{\mathbb{R}}
\newcommand{\C}{\mathbb{C}}
\newcommand{\F}{\mathbb{F}}
\newcommand{\A}{\mathbb{A}}
\newcommand{\PP}{\mathbb{P}}
\newcommand{\E}{\mathbb{E}}
\newcommand{\eps}{\varepsilon}
\newcommand{\dd}{\,\mathrm{d}}
\begin{document}
\section*{UnsolvedMath ID: 700848}
\textbf{Source problem code:} catalogue entry
\subsection{Definitions and mathematical statement}
A group \(G\) is locally finite if every subgroup generated by finitely many elements of \(G\) is finite. Its complex group algebra is
\[
\C[G]=\left\{\sum_{g\in G}a_g g : a_g\in\C,\quad a_g=0\text{ for all but finitely many }g\right\},
\]
with multiplication induced by the group operation and extended bilinearly. The augmentation homomorphism and augmentation ideal are
\[
\varepsilon:\C[G]\longrightarrow\C,\qquad
\varepsilon\left(\sum_{g\in G}a_g g\right)=\sum_{g\in G}a_g,
\qquad I(G)=\ker\varepsilon.
\]
Regard \(I(G)\) as a ring with its inherited addition and multiplication. Here a simple ring means a nonzero ring whose only two-sided ideals are zero and the whole ring; a multiplicative identity is not required. A two-sided ideal means an additive subgroup closed under multiplication on either side by every ring element.

\textbf{Classification problem.} Classify, up to group isomorphism, all locally finite groups \(G\) such that \(I(G)\) is a simple ring. Equivalently, determine precisely which locally finite groups satisfy
\[
I(G)\ne\{0\},\qquad
\forall J\triangleleft I(G),\quad J=\{0\}\ \text{or}\ J=I(G),
\]
where \(J\triangleleft I(G)\) denotes a two-sided ring ideal.

The source does not explicitly specify its convention for simplicity of rings without an identity or for the zero ring. The convention above excludes the trivial group, whose augmentation ideal is zero; this convention requires confirmation against the intended source formulation.
\subsection{Short English statement}
Classify the groups in which every finitely generated subgroup is finite and the augmentation ideal of the complex group algebra is a nonzero simple ring, allowing rings without an identity. The supplied registry labels this problem open; that status has not been independently verified.
\subsection{Sources}
\begin{itemize}
\item Local UnsolvedMath registry, permanent ID \(700848\), exact source problem code \texttt{catalogue entry}, title \emph{Kourovka Notebook Problem 12.29}. Supplied upstream revision: \texttt{d371ee917ac9647fbb06cd6fc529deec6f3dda82}. The record was updated on 2026-10-16 and supplies the status \emph{open}; this is registry metadata, not an independently verified literature assessment.
\item \emph{The Kourovka Notebook (Unsolved Problems in Group Theory)}, Problem 12.29, 12th issue (1992); the supplied background also identifies the 21st issue edition (2026), arXiv:1401.0300v48. The collection metadata names E. I. Khukhro and V. D. Mazurov as editors. The proposer field begins A. E. Zalesski but ends in a malformed diacritic sequence, so its exact intended spelling remains unresolved.
\item The supplied statement attributes the origin of the problem to I. Kaplansky (1965), without a separate bibliographic reference for that attribution.
\item K. Bonvallet, B. Hartley, D. S. Passman, M. K. Smith, \emph{Proc. Amer. Math. Soc.} \textbf{56}, no. 1 (1976), 79--82; supplied as a reference for partial results.
\item A. E. Zalesskiĭ, \emph{Algebra i Analiz} \textbf{2}, no. 6 (1990), 132--149 (Russian); supplied as a reference for partial results.
\item Ch. Praeger, A. E. Zalesskii, \emph{Proc. London Math. Soc.} \textbf{70}, no. 2 (1995), 313--335; supplied as a reference for partial results.
\item B. Hartley, A. Zalesski, \emph{J. London Math. Soc.} \textbf{55} (1997), 210--230; supplied as a reference for partial results.
\end{itemize}
\end{document}
